\section*{Capítulo \ref{ch:b3:bioinorganic} --- Química bioinorgánica}

\begin{solution}{exo:b3:bioinorganic:1}
$\log(\theta/(1 - \theta))$ pasa de $\log 0.25 = -0.602$ a $\log 4 = +0.602$ mientras $\log p$ aumenta en
$\log(5.9/2.1) = 0.449$: $n = 1.204/0.449 = 2.7$.
\end{solution}

\begin{solution}{exo:b3:bioinorganic:2}
$\ce{Mn^2+} < \ce{Fe^2+} < \ce{Co^2+} < \ce{Ni^2+} < \ce{Cu^2+} > \ce{Zn^2+}$. El cinc ($d^{10}$) no tiene estabilización del campo de ligandos ni
distorsión de Jahn--Teller: vuelve a caer por debajo del cobre.
\end{solution}

\begin{solution}{exo:b3:bioinorganic:3}
$\ce{Ca^2+}$ y $\ce{Fe^3+}$, duros: carboxilatos (y fenolatos para el hierro). $\ce{Cu+}$ y $\ce{Hg^2+}$, blandos:
tiolatos de cisteína. $\ce{Zn^2+}$, intermedio: histidina y cisteína.
\end{solution}

\begin{solution}{exo:b3:bioinorganic:4}
$\ce{2H2O -> O2 + 4H+ + 4e-}$: cuatro electrones y cuatro fotones por $\ce{O2}$; cuatro moles de fotones por
mol de $\ce{O2}$.
\end{solution}

\begin{solution}{exo:b3:bioinorganic:5}
Mioglobina: $\theta = 13/13.37 = 0.972$ y $5/5.37 = 0.931$: cede el $(0.972 - 0.931)/0.972 = 4$\,\% de su
oxígeno. Hemoglobina: $\theta = 0.972$ y $0.724$: cede el 26\,\%.
\end{solution}

\begin{solution}{exo:b3:bioinorganic:6}
$[\mathrm{HbCO}]/[\mathrm{HbO_2}] = 230 \times \num{50e-6}/0.2095 = 0.055$; fracción: $0.055/1.055 = 5.2$\,\%.
\end{solution}

\begin{solution}{exo:b3:bioinorganic:7}
Sulfuros, $-2$ cada uno; cisteinatos, $-1$ cada uno. $\ce{[Fe2S2(SR)4]^2-}$: los hierros suman $-2 + 4 + 4 = +6$: dos
Fe(III). $\ce{[Fe2S2(SR)4]^3-}$: $+5$, un Fe(III) y un Fe(II) (valencia mixta). $\ce{[Fe4S4(SR)4]^2-}$:
$-2 + 8 + 4 = +10$, dos Fe(III) y dos Fe(II), cada hierro $+2.5$ en promedio.
\end{solution}

\begin{solution}{exo:b3:bioinorganic:8}
$\ce{[PtCl2(NH3)2] + H2O -> [PtCl(H2O)(NH3)2]+ + Cl-}$. La alta concentración de cloruro de la sangre
desplaza el equilibrio hacia la izquierda; dentro de las células, el cloruro es mucho más bajo y se forma el
acuocomplejo. El platino se une al N7 de la guanina. En el isómero \textit{trans}, las dos
posiciones \omterm{def:b3:complex-mechanisms:labile}{lábiles} están opuestas y no pueden alcanzar dos bases adyacentes de una misma
hebra.
\end{solution}

\begin{solution}{exo:b3:bioinorganic:9}
$1/T_1 = 1/\qty{1.2}{s} + \qty{4.0}{L.mmol^{-1}.s^{-1}} \times \qty{0.10}{mmol/L} = 0.833 + 0.400 = \qty{1.23}{s^{-1}}$: $T_1 = \qty{0.81}{s}$.
\end{solution}

\begin{solution}{exo:b3:bioinorganic:10}
$K = K_{\mathrm{PbL}}/K_{\mathrm{CaL}} = 10^{18.0 - 10.7} = 10^{7.3} = \num{2e7}$: el plomo desplaza por completo al calcio. El ligando
libre uniría también el calcio de la sangre y rebajaría su concentración
peligrosamente; el complejo cálcico se intercambia con el plomo y deja sin cambios el nivel
de calcio.
\end{solution}

\begin{solution}{exo:b3:bioinorganic:11}
Muy por debajo de $K_M$, la reacción es de primer orden: $k' = (k_{\mathrm{cat}}/K_M)[\mathrm E]$, $t_{1/2} = \ln 2/k'$. Anhidrasa
carbónica: $k' = \qty{100}{s^{-1}}$, $t_{1/2} = \qty{6.9}{ms}$. \omterm{def:b3:complex-kinetics:enzyme}{Enzima} típica: $k' = \qty{0.10}{s^{-1}}$, $t_{1/2} = \qty{6.9}{s}$, mil
veces más, demasiado lento para el segundo que pasa un glóbulo rojo en un capilar pulmonar.
\end{solution}

\begin{solution}{exo:b3:bioinorganic:12}
La mitad de los sitios con CO significa $[\mathrm{HbCO}] = [\mathrm{HbO_2}]$: $p(\ce{CO}) = p(\ce{O2})/M = 0.2095/230 = \qty{9.1e-4}{bar}$, unas
\qty{910}{ppm}.
\end{solution}

\begin{solution}{pb:b3:bioinorganic:1}
\textbf{1.} $\theta = 13^{2.7}/(3.5^{2.7} + 13^{2.7}) = 0.972$.

\textbf{2.} $\theta(\qty{5.0}{kPa}) = 0.724$.

\textbf{3.} $5.0/(0.37 + 5.0) = 0.931$.

\textbf{4.} Su $p_{50}$ es diez veces menor: a las presiones de un músculo activo está mucho más
saturada que la hemoglobina a la misma presión, de modo que el oxígeno pasa de la
hemoglobina a la mioglobina.

\textbf{5.} $n = 1$.

\textbf{6.} La unión en un \omterm{def:b3:bioinorganic:haem}{hemo} desplaza su hierro hacia el plano y arrastra la histidina y
su hélice; las subunidades cambian juntas a una forma de mayor afinidad por los
demás sitios.

\textbf{7.} $\qty{150}{g/L}/\qty{64500}{g/mol} = \qty{2.33e-3}{mol/L}$.

\textbf{8.} $4 \times \num{2.33e-3} = \qty{9.30e-3}{mol/L}$.

\textbf{9.} $\qty{9.30}{mmol/L} \times 0.972 = \qty{9.04}{mmol}$.

\textbf{10.} $\qty{9.30}{mmol/L} \times (0.972 - 0.724) = \qty{2.31}{mmol}$.

\textbf{11.} $\qty{2.31e-3}{mol} \times \qty{22.41}{L/mol} = \qty{52}{mL}$.

\textbf{12.} $0.248/0.972 = 26$\,\%.

\textbf{13.} $5.0^{2.7}/(4.0^{2.7} + 5.0^{2.7}) = 0.646$.

\textbf{14.} $\qty{9.30}{mmol/L} \times (0.972 - 0.646) = \qty{3.03}{mmol}$.

\textbf{15.} $3.03/2.31 = 1.31$: un treinta por ciento más.

\textbf{16.} El dióxido de carbono producido por la respiración, hidratado a ácido carbónico por la
anhidrasa carbónica, y el ácido láctico en el esfuerzo intenso.

\textbf{17.} $\qty{3.03}{mmol/L} \times \qty{5.0}{L/min} = \qty{15.1}{mmol/min}$.

\textbf{18.} $\qty{15.1e-3}{mol/min} \times \qty{22.41}{L/mol} = \qty{0.34}{L/min}$.

\textbf{19.} $230 \times \num{100e-6}/0.2095 = 0.110$.

\textbf{20.} $0.110/1.110 = 0.099$: alrededor del 10\,\% de los sitios.

\textbf{21.} Los sitios que aún llevan oxígeno, en una molécula que lleva también CO, están en
la forma de alta afinidad: la curva se desplaza hacia la izquierda y ceden menos de
su oxígeno en los tejidos.

\textbf{22.} Según la relación de competencia, la razón $[\mathrm{HbCO}]/[\mathrm{HbO_2}]$ disminuye en proporción a $p(\ce{O2})$:
respirar oxígeno puro, a casi cinco veces la presión que tiene en el aire, desplaza el monóxido de
carbono más deprisa.

\textbf{23.} Al dificultar la unión lineal del CO y formar un enlace de hidrógeno con el $\ce{O2}$, mantiene $M$ en
los cientos en lugar del valor mucho mayor del \omterm{def:b3:bioinorganic:haem}{hemo} libre; de lo contrario, incluso el
monóxido de carbono producido en el organismo bloquearía buena parte de la hemoglobina.

\textbf{24.} En reposo, la sangre entrega unos \qty{0.34}{L} de oxígeno por minuto.
\end{solution}
